\documentclass[12pt]{article}
\usepackage[utf8]{inputenc}
\usepackage[T1]{fontenc}
\usepackage{hyperref}
\title{{\~N}andeva Language Resources for K-12 Curriculum Development}
\author{Compiled by Claude Sonnet 5 (claude-sonnet-5) for Andr\'as Kornai}
\date{}
\begin{document}
\maketitle
\section{Introduction}
{\~N}andeva (ISO 639-3 tpj; Glottocode tapi1253) is a Tupian language of Paraguay, Argentina and Bolivia. It is endangered. Names overlap across sources: ``{\~N}andeva'' is also used for Chirip\'a and for Ava Guaran\'i, and the tpj code is linked to the Tapiete. Any resource must be checked for which variety it covers. Searches found no dictionary, primer or school materials specific to tpj.
\section{Primary Educational Resources}
None found in search results. The language is not known to be taught in schools.
\section{Linguistic and Reference Materials}
\begin{itemize}
\item Wikipedia, Guarani languages: \url{https://en.wikipedia.org/wiki/Guarani_languages}
\item Wikipedia, Guarani dialects: \url{https://en.wikipedia.org/wiki/Guarani_dialects}
\end{itemize}
\section{Religious and Community Resources}
The Native Languages page on Chirip\'a lists Bible translations and prayers under ``Words of Life: Chiripa''. That is a different variety from tpj: \url{https://www.native-languages.org/chiripa.htm}
\section{Digital and Audio Resources}
\begin{itemize}
\item Endangered Languages Project entry (nhd), which has metadata and an OLAC search link: \url{https://endangeredlanguages.com/lang/2920}
\item Native Land Digital page on Guarani, which links Paraguayan sources including the Secretar\'ia de Pol\'iticas Ling\"u\'isticas: \url{https://native-land.ca/maps/languages/guarani-2}
\end{itemize}
\section{Curriculum Development Context}
Little exists. Paraguayan Guaran\'i materials are abundant but cover a different variety. Any curriculum would need documentation work with the communities first.
\end{document}